\begin{table}[H]
    \centering
    \caption{Dampak Skenario SLODR terhadap Ekor Distribusi Skor Komposit (Marginal Johnson SB, \textit{Copula} Gaussian)}
    \label{tab:slodr_ekor}
    \begin{spacing}{1.0}
    \resizebox{\textwidth}{!}{%
    \begin{tabular}{lcccccc}
        \toprule
        & Normal & \multicolumn{5}{c}{Kemiringan SLODR $b$} \\
        \cmidrule(lr){3-7}
        & & 0,00 & 0,01 & 0,02 & 0,03 & 0,04 \\
        \midrule
        Skewness komposit & $0{,}000$ & 0,419 & 0,361 & 0,304 & 0,249 & 0,195 \\
        \midrule
        \multicolumn{7}{l}{\textit{Peluang ekor (\%)}} \\
        $M - 3\,S_Z$ & 0,135 & $<0{,}001$ & $<0{,}001$ & $<0{,}001$ & 0,001 & 0,003 \\
        $M - 2\,S_Z$ & 2,275 & 0,812 & 0,965 & 1,110 & 1,273 & 1,432 \\
        $M + 2\,S_Z$ & 2,275 & 3,327 & 3,180 & 3,026 & 2,881 & 2,716 \\
        $M + 3\,S_Z$ & 0,135 & 0,380 & 0,318 & 0,257 & 0,206 & 0,161 \\
        \midrule
        \multicolumn{7}{l}{\textit{Skor komposit pada persentil normal}} \\
        $M - 3\,S_Z$ & 296,9 & 341,8 & 338,4 & 335,3 & 332,1 & 328,9 \\
        $M - 2\,S_Z$ & 364,6 & 381,7 & 379,8 & 377,8 & 375,9 & 374,1 \\
        $M + 2\,S_Z$ & 635,4 & 648,9 & 647,0 & 644,8 & 642,6 & 640,7 \\
        $M + 3\,S_Z$ & 703,1 & 730,3 & 724,4 & 718,4 & 712,3 & 707,1 \\
        \bottomrule
    \end{tabular}%
    }
    \end{spacing}
    \vspace{2mm}
    \begin{flushleft}
    \footnotesize\textit{Catatan.} SLODR adalah \textit{Spearman's Law of Diminishing Returns}, yaitu gagasan bahwa korelasi antartes menurun pada kemampuan tinggi. Muatan \textit{g-factor} ditetapkan sebagai $m(g) = m_0 - b\,g$ dengan $g \sim N(0,1)$; $b = 0$ berarti korelasi konstan. $m_0$ dikalibrasi pada tiap $b$ sehingga $SD$ komposit $= S_Z = 67,72$; skor komposit dipusatkan ke $500$. Peluang ekor mengacu pada sisi bawah untuk $k<0$ dan sisi atas untuk $k>0$. Korelasi bersyarat pada $g=-1$, $+1$, dan $+2$ masing-masing $0,446$, $0,394$, dan $0,370$ pada $b = 0,02$.
    \end{flushleft}
\end{table}
